\chapter{只剩两个载体位置时，直接指定最后的余数}\label{ch:two}
\section{不要为所有输入预先选择同一条路}
现在 $n=2p+2$。核心写成
\[((u+v)^p,(u-v)^p,(-pu)^2).
\]
令 $m=p+1$，合法本原参数要求 $\gcd(u,v)=1$ 和 $\gcd(v,m)=1$，所以 $v$ 是奇数。

这一讲采用更有针对性的办法。先取得一些剪切和一个足够深的主同余子群；再根据当前输入选两次返回，直接把第二坐标设为 $1$ 或 $-1$ 模控制层。这样不必先证明所有有限尺度都可达。
\InsertProof{prop:r2}

\section{两次CRT分别负责什么}
已有上剪切就能把 $u$ 改成 $u+kv$。由于 $v$ 奇，可以把它写成 $2a-v$。某个数字返回给
\[(u,v)=(2a-v,v)quad\longmapsto\quad(-a,pv-ma)
\]
的方向。现在 $a$ 可以在一个等差数列中自由选，我们用它做两件事。

第一次，要让 $pv-ma$ 在控制模数 $q$ 的每个素因子处都非零。每个素数只提出一个条件，用 CRT 合并。不是让这个数变小，而是让它成为模 $q$ 的单位。

第二次，既然它已经可逆，就能选一个允许的数字 $b$，再选新的平移参数，让
\[bv_1-ma_1\equiv\epsilon\pmod q,qquad\epsilon=\pm1.
\]
与此同时还要让输出坐标互素；否则最后约分可能把刚安排的余数改掉。证明中的额外CRT条件正是在保护本原性。

\section{一组可以手算的CRT数据}
取 $p=5,m=6$，该构造可使用 $q=(p^2-1)^2=576$。从合法参数 $(u,v)=(1,5)$ 开始，先选 $a=8$，满足 $2a\equiv1\pmod5$。上平移把第一坐标变成 $2a-v=11$，再用第一个数字返回，得到本原分析列
\[(-8,25-48)=(-8,-23).
\]
$23$ 与 $576$ 互素，第一次任务完成。

$-23\equiv1\pmod6$。为避免不可用的数字 $b=1$，取 $b=5,\epsilon=-1$。第二个平移参数要满足
\[a_1\equiv-4\pmod{23},qquad a_1\equiv-19\pmod{96}.
\]
可取 $a_1=1997$，它还与五互素。第二次输出为
\[(-1997,5(-23)-6\cdot1997)=(-1997,-12097).
\]
两坐标最大公因数是一，且 $-12097\equiv-1\pmod{576}$。再用上剪切 $U(307)$，第一坐标成为 $-3715776=-6451\cdot576$。这就准确送进了 $(0,-1)$ 的同余类，精确输送引理接手。

数值在这个分析代表里变得很大，但实际平均词带着共同收缩因子。这个例子强调的是\textbf{有一个有限的解法}；它没有声称这条解法短，也没有声称每次本原参数都变小。

\section{把这些数字变成真实位置选择}
下面先给上平移和缩放逆的平均词，再给横向反射，最后证明两次CRT可对所有合法输入执行。刚才的例子对应证明的第三、第四阶段。
\InsertProof{lem:two-carrier}
